% -*- mode: latex; coding: utf-8 -*-
%-----------------------------------------------------------------------------
% pxDIAG -- User manual and API reference
% Part of P3CTeX. Author: Scvria.
%
% Build from directory tex/ (with TEXINPUTS=./latex;./code;./doc;):
%   pdflatex -interaction=nonstopmode -output-directory=doc doc/pxDIAG.tex
%   pdflatex -interaction=nonstopmode -output-directory=doc doc/pxDIAG.tex
%-----------------------------------------------------------------------------
\documentclass[11pt,a4paper]{article}
\PassOptionsToPackage{colorlinks,urlcolor=blue,linkcolor=blue}{hyperref}
\usepackage[T1]{fontenc}
\usepackage[utf8]{inputenc}
\usepackage[margin=22mm]{geometry}
\usepackage{lmodern}
\usepackage[catalan]{babel}
\usepackage[preset=p3c-uoc]{pxDIAG}
\usepackage{hyperref}

\title{El paquet \texttt{pxDIAG}\\
  \large Diagrames de Tombstone, compiladors, aut\`omats i l\`ogica computacional per a \LaTeX2e}
\author{Scvria\\P3CTeX / UOC PEC Repository}
\date{2026/10/27 -- Versi\'o 0.1}

\begin{document}
\maketitle

\begin{abstract}
  \texttt{pxDIAG} proporciona una interf\'icie \LaTeX2e\ ergon\`omica basada en llistes i claus
  (\texttt{l3keys}) per construir diagrames de compiladors i computaci\'o dins de l'ecosistema
  \textbf{P3CTeX} (o de forma independent en classes est\`andard com \texttt{article}).
  Cobreix cinc grans fam\'ilies: (1) diagrames de Tombstone (diagrames en T),
  (2) diagrames de blocs de traductors, pipelines i codi intermedi,
  (3) arbres sint\`actics, aut\`omats finits i grafs de computaci\'o,
  (4) gram\`atiques BNF, derivacions, taules l\`exiques i comparatives de codi de tres adreces,
  i (5) regles d'infer\`encia l\`ogica, diagrames de Venn/jerarquies i diagrames de flux.
\end{abstract}

\tableofcontents
\newpage

\section{C\`arrega i activaci\'o del paquet}

Dins d'un document amb la classe \texttt{P3CTeX}, activeu el m\`odul opcional \texttt{CMP}
a les opcions de classe (sense perdre els m\`oduls per defecte):

\begin{verbatim}
\documentclass[diagrams,PAC,default]{P3CTeX}
\end{verbatim}

Tamb\'e es pot carregar a trav\'es de \texttt{pxCORE} o de manera independent en qualsevol
document est\`andard:

\begin{verbatim}
% A traves de pxCORE amb opcions:
\usepackage[diagrams,diagrams-config={preset=p3c-blueprint}]{pxCORE}

% O directament de manera independent:
\usepackage[preset=p3c-uoc]{pxDIAG}
\end{verbatim}

\section{Configuraci\'o i Presets (\texttt{pxDIAG})}

 Totes les opcions es poden ajustar globalment (o dins d'un grup) amb \verb|\pxDiagramSetup{...}|
o localment al primer argument opcional \verb|[...]| de qualsevol ordre.

\subsection{Presets predefinits}
\begin{itemize}
  \item \texttt{p3c-uoc} (per defecte): Estil inspirat en els m\`oduls did\`actics de la UOC
    (\texttt{PID\_00215089}), amb peces blanques, ombra tridimensional i panell de fons suau.
  \item \texttt{p3c-blueprint}: Paleta corporativa blava/cian de P3CTeX.
  \item \texttt{p3c-minimal}: Disseny net sense marc exterior ni ombres.
  \item \texttt{p3c-exam}: Alt contrast per a impressi\'o i resoluci\'o d'ex\`amens o PACs.
\end{itemize}

\begin{verbatim}
\pxDiagramUsePreset{p3c-blueprint}
\pxDiagramSavePreset{el-meu-estil}{framed=true, shadow=false, scale=0.9}
\end{verbatim}

\subsection{Claus principals}
\begin{itemize}
  \item \texttt{unit} (per defecte \texttt{1.18cm}): Unitat modular de la graella de Tombstone.
  \item \texttt{scale} (per defecte \texttt{1}): Factor d'escala del diagrama TikZ.
  \item \texttt{width} / \texttt{max-width} (per defecte \texttt{\textbackslash linewidth}): Redimensionament autom\`atic.
  \item \texttt{framed} (\texttt{true}/\texttt{false}): Panell exterior estil UOC.
  \item \texttt{shadow} (\texttt{true}/\texttt{false}): Ombra de les peces de Tombstone.
  \item \texttt{title}, \texttt{caption}, \texttt{label}, \texttt{float}: T\'itol del panell i suport per a figures flotants.
\end{itemize}

\newpage
\section{Fam\'ilia 1: Diagrames de Tombstone (Diagrames en T)}

\subsection{Peces b\`asiques}
Les quatre peces fonamentals dels diagrames de Tombstone es poden invocar de manera
aut\`onoma o col$\cdot$locar-se en coordenades \texttt{[x,y]} dins d'un entorn \texttt{pxDiagram}:

\begin{verbatim}
\pxTCompiler{Pascal}{EXE-PC}{EXE-PC}
\pxTProgram{ordenar}{Pascal}
\pxTMachine{LINUX-PC}
\pxTInterpreter{Basic}{EXE-PC}
\end{verbatim}

\pxTCompiler[inline=true]{Pascal}{EXE-PC}{EXE-PC}\quad
\pxTProgram[inline=true]{ordenar}{Pascal}\quad
\pxTMachine[inline=true]{LINUX-PC}\quad
\pxTInterpreter[inline=true]{Basic}{EXE-PC}

\subsection{Composicions autom\`atiques (sense coordenades)}

\subsubsection*{Execuci\'o, interpretaci\'o i compilaci\'o d'un programa}
Amb l'opci\'o \texttt{show-rules=true} es mostren les fletxes de coher\`encia (\emph{Iguals})
en els punts d'uni\'o:

\begin{verbatim}
\pxTombstoneRun[show-rules=true]{P}{M}
\pxTombstoneInterpret[show-rules=true]{P}{L}{M}
\pxTombstoneCompile[show-rules=true]{P}{L1}{L2}{M}
\end{verbatim}

\pxTombstoneRun[show-rules=true,title={Execuci\'o d'un programa}]{P}{M}
\pxTombstoneInterpret[show-rules=true,title={Interpretaci\'o d'un programa}]{P}{L}{M}
\pxTombstoneCompile[show-rules=true,title={Compilaci\'o d'un programa}]{P}{L1}{L2}{M}

\subsubsection*{M\`aquina virtual i compiladors interpretatius (Java / JVM)}
\begin{verbatim}
\pxTombstoneVM[title={M\`aquina virtual}]{L}{M}
\end{verbatim}
\pxTombstoneVM[title={M\`aquina virtual}]{L}{M}

\subsubsection*{Compilador enlla\c{c}ador i cadenes de traducci\'o}
\begin{verbatim}
\pxTombstoneChain[title={Compilador enlla\c{c}ador}]{ordenar}{
  Ada/OBJ/M/Compilador,
  OBJ/M/M/Enlla\c{c}ador
}
\end{verbatim}
\pxTombstoneChain[title={Compilador enlla\c{c}ador}]{ordenar}{
  Ada/OBJ/M/Compilador,
  OBJ/M/M/Enlla\c{c}ador
}

\subsubsection*{Uni\'o de dos compiladors, compilaci\'o creuada i Bootstrapping}
\begin{verbatim}
\pxTombstoneCompileCompiler[id=(1),id2=(2),show-bindings=true]
  {Pascal}{Pentium}{C}{C}{Pentium}{Pentium}
\pxTombstoneBootstrap[title={Compilaci\'o creuada / Bootstrapping}]{Pascal}{M4}{M5}
\end{verbatim}

\pxTombstoneCompileCompiler[id=(1),id2=(2),show-bindings=true,title={Uni\'o de dos compiladors}]{Pascal}{Pentium}{C}{C}{Pentium}{Pentium}
\pxTombstoneBootstrap[title={Compilaci\'o creuada / Bootstrapping en dues fases}]{Pascal}{M4}{M5}

\subsection{Entorn de llen\c{c} componible (\texttt{pxDiagram})}
Per a exercicis complets estil UOC (\emph{Tenim / Volem / S'ha de programar / Proc\'es}),
l'entorn \texttt{pxDiagram} permet combinar peces i notes en una graella exacta on
dos compiladors encaixen desplacant $\Delta x = 2.0$, $\Delta y = -0.7$:

\begin{verbatim}
\begin{pxDiagram}[title={Exercici: Compilador de Prolog sobre int\`erpret}]
  \pxDIAGNote[font=\sffamily\bfseries\footnotesize]{-2.8, 2.7}{Tenim:}
  \pxTInterpreter{Prolog}{P-IV}[0, 2.2]
  \pxTMachine{P-IV}[][2.0, 3.1]
  \pxDIAGNote[font=\sffamily\bfseries\footnotesize]{-2.8, 0.35}{Proc\'es:}
  \pxTCompiler{Prolog}{P-IV}{Prolog}[][0.5, 0]
  \pxTCompiler{Prolog}{P-IV}{Prolog}[][2.5, -0.7]
  \pxTInterpreter{Prolog}{P-IV}[2.5, -2.3]
  \pxTMachine{P-IV}[][2.5, -2.3]
  \pxTCompiler{Prolog}{P-IV}{P-IV}[][4.5, 0]
\end{pxDiagram}
\end{verbatim}

\begin{pxDiagram}[title={Exercici: Compilador de Prolog sobre int\`erpret}]
  \pxDIAGNote[font=\sffamily\bfseries\footnotesize]{-2.8, 2.7}{Tenim:}
  \pxTInterpreter{Prolog}{P-IV}[0, 2.2]
  \pxTMachine{P-IV}[][2.0, 3.1]
  \pxDIAGNote[font=\sffamily\bfseries\footnotesize]{-2.8, 0.35}{Proc\'es:}
  \pxTCompiler{Prolog}{P-IV}{Prolog}[][0.5, 0]
  \pxTCompiler{Prolog}{P-IV}{Prolog}[][2.5, -0.7]
  \pxTInterpreter{Prolog}{P-IV}[2.5, -2.3]
  \pxTMachine{P-IV}[][2.5, -2.3]
  \pxTCompiler{Prolog}{P-IV}{P-IV}[][4.5, 0]
\end{pxDiagram}

\newpage
\section{Fam\'ilia 2: Diagrames de Blocs de Traductors i Pipelines}

\begin{verbatim}
\pxTranslatorBox{Programa\\font}{Traductor}{Programa\\objecte}[Errors]
\pxCompileRunFlow{Programa\\font}{Compilador}{Programa\\objecte}
  {Plataforma\\d'execuci\'o}{Resultats}[Errors][Dades]
\pxCompilerModel[title={Model d'un compilador}]
\pxIRHub{Cobol, C, Ada, Modula, Pascal}{Codi\\intermedi}
  {80x86, Mac, Pentium IV, Linux-PC, Aix-IBM}
\end{verbatim}

\pxTranslatorBox[title={Funcionament d'un traductor}]{Programa\\font}{Traductor}{Programa\\objecte}[Errors]
\pxCompileRunFlow[title={Proc\'es de compilaci\'o i execuci\'o}]{Programa\\font}{Compilador}{Programa\\objecte}{Plataforma\\d'execuci\'o}{Resultats}[Errors][Dades]
\pxCompilerModel[title={Model d'un compilador}]
\pxIRHub[title={Estrat\`egia d'\'us del codi intermedi}]{Cobol, C, Ada, Modula, Pascal}{Codi\\intermedi}{80x86, Mac, Pentium IV, Linux-PC, Aix-IBM}

\newpage
\section{Fam\'ilia 3: Arbres Sint\`actics, Aut\`omats i Grafs}

\subsection{Arbres sint\`actics (\texttt{\textbackslash pxSyntaxTree})}
Els nodes amb tres camps (\texttt{id / etiqueta / x,y}) es dibuixen com a variables sint\`actiques
no terminals, mentre que els nodes amb quatre camps (\texttt{id / TESTIMONI / x,y / lexema})
generen autom\`aticament la doble caixa de testimoni i valor l\`exic:

\begin{verbatim}
\pxSyntaxTree[title={Arbre sint\`actic}]{
  exp / exp / 0,0;
  assigna / assigna / -2.2,-1.1;
  semi / PUNT\_I\_COMA / 2.2,-1.1 / {;};
  id1 / IDENTIFICADOR / -4.4,-2.4 / ValorX;
  eq / ASSIGNACI\'O / -2.2,-2.4 / {=};
  op / operaci\'o / 0.2,-2.4;
  opd1 / operand / -1.8,-3.7;
  opr / operador / 0.2,-3.7;
  opd2 / operand / 2.2,-3.7;
  id2 / IDENTIFICADOR / -1.8,-4.9 / ValorY;
  plus / SUMA / 0.2,-4.9 / {+};
  num / ENTER / 2.2,-4.9 / 1
}{
  exp->assigna, exp->semi, assigna->id1, assigna->eq, assigna->op,
  op->opd1, op->opr, op->opd2, opd1->id2, opr->plus, opd2->num
}
\end{verbatim}

\pxSyntaxTree[title={Arbre sint\`actic}]{
  exp / exp / 0,0;
  assigna / assigna / -2.2,-1.1;
  semi / PUNT\_I\_COMA / 2.2,-1.1 / {;};
  id1 / IDENTIFICADOR / -4.4,-2.4 / ValorX;
  eq / ASSIGNACI\'O / -2.2,-2.4 / {=};
  op / operaci\'o / 0.2,-2.4;
  opd1 / operand / -1.8,-3.7;
  opr / operador / 0.2,-3.7;
  opd2 / operand / 2.2,-3.7;
  id2 / IDENTIFICADOR / -1.8,-4.9 / ValorY;
  plus / SUMA / 0.2,-4.9 / {+};
  num / ENTER / 2.2,-4.9 / 1
}{
  exp->assigna, exp->semi, assigna->id1, assigna->eq, assigna->op,
  op->opd1, op->opr, op->opd2, opd1->id2, opr->plus, opd2->num
}

\subsection{Aut\`omats finits (\texttt{\textbackslash pxAutomaton}) i Grafs (\texttt{\textbackslash pxGraph})}
\begin{verbatim}
\pxAutomaton[title={Aut\`omat finit determinista}]{
  q0 / $q_0$ / 0,0 / initial;
  q1 / $q_1$ / 3.2,0 / accepting
}{
  q0 -> q1 / lletra;
  q1 -> q1 / {lletra, d\'igit} / loop above
}
\end{verbatim}

\pxAutomaton[title={Aut\`omat finit determinista}]{
  q0 / $q_0$ / 0,0 / initial;
  q1 / $q_1$ / 3.2,0 / accepting
}{
  q0 -> q1 / lletra;
  q1 -> q1 / {lletra, d\'igit} / loop above
}

\newpage
\section{Fam\'ilia 4: Gram\`atiques BNF, Derivacions, Testimonis i Codi Intermedi}

\begin{verbatim}
\pxLexemeTable[title={An\`alisi l\`exica}][ValorX = ValorY + 1;]{
  IDENTIFICADOR / ValorX; ASSIGNACI\'O / {=}; IDENTIFICADOR / ValorY;
  SUMA / {+}; ENTER / 1; PUNT\_I\_COMA / {;}
}
\pxBNF[title={Gram\`atica lliure de context}]{
  exp -> assigna PUNT\_I\_COMA;
  assigna -> IDENTIFICADOR ASSIGNACI\'O operaci\'o;
  operaci\'o -> operand operador operand
}
\pxIRCompare[title={Optimitzaci\'o sobre codi de tres direccions}]
  {Codi de tres direccions}
  {E1: R1 := A + 3; IF R1 > B THEN GO TO E2; E2: A := A + 1}
  {Codi optimitzat}
  {E1: R := A + 3; IF R <= 10 THEN GO TO E3; E2: INC (A)}
\end{verbatim}

\pxLexemeTable[title={An\`alisi l\`exica}][ValorX = ValorY + 1;]{
  IDENTIFICADOR / ValorX;
  ASSIGNACI\'O / {=};
  IDENTIFICADOR / ValorY;
  SUMA / {+};
  ENTER / 1;
  PUNT\_I\_COMA / {;}
}

\pxBNF[title={Gram\`atica lliure de context}]{
  exp -> assigna PUNT\_I\_COMA;
  assigna -> IDENTIFICADOR ASSIGNACI\'O operaci\'o;
  operaci\'o -> operand operador operand
}

\pxIRCompare[title={Optimitzaci\'o sobre codi de tres direccions}]
  {Codi de tres direccions}
  {E1: R1 := A + 3; IF R1 > B THEN GO TO E2; E2: A := A + 1}
  {Codi optimitzat}
  {E1: R := A + 3; IF R <= 10 THEN GO TO E3; E2: INC (A)}

\section{Fam\'ilia 5: L\`ogica, Conjunts i Diagrames de Flux}

\begin{verbatim}
\pxInferenceRule{\Gamma \vdash e_1 : \tau_1, \Gamma \vdash e_2 : \tau_1}
  {\Gamma \vdash e_1 = e_2 : \texttt{bool}}[T-Eq]
\pxLogicGroup[title={Jerarquia de Chomsky}]{
  Tipus 3 (Regulars), Tipus 2 (Lliures de context),
  Tipus 1 (Sensibles al context), Tipus 0 (Recursivament enumerables)
}
\end{verbatim}

\pxInferenceRule{\Gamma \vdash e_1 : \tau_1, \Gamma \vdash e_2 : \tau_1}{\Gamma \vdash e_1 = e_2 : \texttt{bool}}[T-Eq]

\pxLogicGroup[title={Jerarquia de Chomsky}]{
  Tipus 3 (Regulars), Tipus 2 (Lliures de context),
  Tipus 1 (Sensibles al context), Tipus 0 (Recursivament enumerables)
}

\section{Fam\'ilia 6: Expressions Regulars, Mem\`oria i Endianness (v0.2)}

Per a exercicis d'especificaci\'o l\`exica (JLex/Flex), diagrames de Tombstone amb inc\`ognites ($A, B, C$) i arquitectura de sistemes encastats (SSEE / Estructura de Computadors), \texttt{pxDIAG} incorpora:

\begin{verbatim}
\pxTombstoneCompile[title={Tombstone amb inc\`ognites}]{Prog}{Swift/A}{AMD64}{B}[C]
\pxRegEx{"@"[a-zA-Z][a-zA-Z0-9_]*}[Patr\'o JLex per a noms d'usuari]
\pxEndianCompare[title={Disposici\'o a mem\`oria}]{Paraula de 32 bits: \texttt{0x001C5880}}{
  0x5AF4 / 0x80 (LSB) / 0x00 (MSB);
  0x5AF5 / 0x58       / 0x1C;
  0x5AF6 / 0x1C       / 0x58;
  0x5AF7 / 0x00 (MSB) / 0x80 (LSB)
}
\pxRegisterMap[title={Registre GPIOx\_MODER}]{
  31:16 / MODER15..8 / 0xAAAA;
  15:4  / MODER7..2  / 0x555;
  3:2   / MODER1     / 01;
  1:0   / MODER0     / 01
}
\end{verbatim}

\pxTombstoneCompile[title={Tombstone amb inc\`ognites}]{Prog}{Swift/A}{AMD64}{B}[C]

\pxRegEx{"@"[a-zA-Z][a-zA-Z0-9_]*}[Patr\'o JLex per a noms d'usuari]

\pxEndianCompare[title={Disposici\'o a mem\`oria}]{Paraula de 32 bits: \texttt{0x001C5880}}{
  0x5AF4 / 0x80 (LSB) / 0x00 (MSB);
  0x5AF5 / 0x58       / 0x1C;
  0x5AF6 / 0x1C       / 0x58;
  0x5AF7 / 0x00 (MSB) / 0x80 (LSB)
}

\pxRegisterMap[title={Registre GPIOx\_MODER}]{
  31:16 / MODER15..8 / 0xAAAA;
  15:4  / MODER7..2  / 0x555;
  3:2   / MODER1     / 01;
  1:0   / MODER0     / 01
}

\newpage
\section{Fam\'ilia 7: Grafs de Cerca, Taules d'Aut\`omats, Machine Learning i Mineria de Dades}

Per donar resposta a les assignatures de \emph{Grafs i Complexitat}, \emph{Aut\`omats i Gram\`atiques}, \emph{Intel\textperiodcentered lig\`encia Artificial}, \emph{Aprenentatge Computacional} i \emph{Mineria de Dades}, \texttt{pxDIAG} incorpora:

\subsection{Taules de transici\'o i traces d'aut\`omats (\texttt{\textbackslash pxTransitionTable}, \texttt{\textbackslash pxAutomatonTrace})}

\begin{verbatim}
\pxTransitionTable[title={Taula de transici\'o i partici\'o $\Pi_0$}]{a, b}{
  q_0 / initial   / q_1, q_2 / I;
  q_1 / normal    / q_3, q_2 / I;
  q_2 / accepting / q_1, q_3 / II;
  q_3 / accepting / q_3, q_3 / II
}
\pxAutomatonTrace[title={Tra\c{c}a de la cadena \texttt{aba}}]{
  q_0 / a / q_1; q_1 / b / q_2; q_2 / a / q_1
}[Rebutjat]
\end{verbatim}

\pxTransitionTable[title={Taula de transici\'o i partici\'o $\Pi_0$}]{a, b}{
  q_0 / initial   / q_1, q_2 / I;
  q_1 / normal    / q_3, q_2 / I;
  q_2 / accepting / q_1, q_3 / II;
  q_3 / accepting / q_3, q_3 / II
}

\pxAutomatonTrace[title={Tra\c{c}a de la cadena \texttt{aba}}]{
  q_0 / a / q_1; q_1 / b / q_2; q_2 / a / q_1
}[Rebutjat]

\subsection{Grafs de cerca heur\'istica i camins \`optims (\texttt{\textbackslash pxGraph})}

\begin{verbatim}
\pxGraph[title={Graf de cerca heur\'istica $A^*$}]{
  A / A / 0,0 / start / $h=20$;
  B / B / 2.2,1.1 / path / {$g=4, f=16$};
  C / C / 2.2,-1.1 / visited / $h=15$;
  D / D / 4.4,1.1 / open / $h=10$;
  G / G / 4.4,-1.1 / goal / $h=0$
}{
  A => B / 4; A -> C / 5; B -> D / 6; B => G / 8; C -> G / 12
}
\end{verbatim}

\pxGraph[title={Graf de cerca heur\'istica $A^*$}]{
  A / A / 0,0 / start / $h=20$;
  B / B / 2.2,1.1 / path / {$g=4, f=16$};
  C / C / 2.2,-1.1 / visited / $h=15$;
  D / D / 4.4,1.1 / open / $h=10$;
  G / G / 4.4,-1.1 / goal / $h=0$
}{
  A => B / 4; A -> C / 5; B -> D / 6; B => G / 8; C -> G / 12
}

\subsection{Arbres de decisi\'o i Xarxes Neuronals (\texttt{\textbackslash pxDecisionTree}, \texttt{\textbackslash pxNeuralNet})}

\begin{verbatim}
\pxDecisionTree[title={Arbre de decisi\'o (ID3 / C4.5)}]{
  neu   / Neu   / 0,0       / root     / $IG=0.875$;
  neuL  / {Y\\$100\%$} / -2.8,-1.8 / leaf;
  mvia  / Via   / 0,-1.8    / decision / $IG=0.722$;
  neuH  / {N\\$100\%$} / 2.8,-1.8  / leaf;
  mviaA / {Y\\$100\%$} / -1.4,-3.6 / leaf;
  mviaC / {N\\$100\%$} / 1.4,-3.6  / leaf
}{
  neu -> neuL / L; neu -> mvia / M; neu -> neuH / H;
  mvia -> mviaA / A; mvia -> mviaC / C
}

\pxNeuralNet[title={Perceptr\'o Multicapa (MLP)}]{
  Entrada / x_0, x_1, x_2 / input;
  Capa oculta / h_0, h_1, h_2, h_3 / ReLU;
  Sortida / Y / Heaviside
}[x_0 -> h_0 / 1; x_1 -> h_0 / -1; h_0 -> Y / -1; h_2 -> Y / 1]
\end{verbatim}

\pxDecisionTree[title={Arbre de decisi\'o (ID3 / C4.5)}]{
  neu   / Neu   / 0,0       / root     / $IG=0.875$;
  neuL  / {Y\\$100\%$} / -2.8,-1.8 / leaf;
  mvia  / Via   / 0,-1.8    / decision / $IG=0.722$;
  neuH  / {N\\$100\%$} / 2.8,-1.8  / leaf;
  mviaA / {Y\\$100\%$} / -1.4,-3.6 / leaf;
  mviaC / {N\\$100\%$} / 1.4,-3.6  / leaf
}{
  neu -> neuL / L; neu -> mvia / M; neu -> neuH / H;
  mvia -> mviaA / A; mvia -> mviaC / C
}

\pxNeuralNet[title={Perceptr\'o Multicapa (MLP)}]{
  Entrada / x_0, x_1, x_2 / input;
  Capa oculta / h_0, h_1, h_2, h_3 / ReLU;
  Sortida / Y / Heaviside
}[x_0 -> h_0 / 1; x_1 -> h_0 / -1; h_0 -> Y / -1; h_2 -> Y / 1]

\subsection{Matrius de confusi\'o i L\`ogica borrosa (\texttt{\textbackslash pxConfusionMatrix}, \texttt{\textbackslash pxFuzzyPlot})}

\begin{verbatim}
\pxConfusionMatrix[title={Matriu de confusi\'o i m\`etriques F1}]
  {Positiu (1), Negatiu (0)}
  {3 (TP), 2 (FP); 5 (FN), 5 (TN)}
  [Precisi\'o = 0.60 \quad|\quad Sensibilitat = 0.375 \quad|\quad $F_1 = 0.4615$]

\pxFuzzyPlot[title={Funcions de pertinen\c{c}a borrosa}, xlabel={Temperatura ($^\circ$C)}]{0, 50}{
  Baixa (VL)  / primaryColor           / 0,1 / 10,1 / 20,0;
  Mitjana (M) / secondaryColor!75!black / 12,0 / 22,1 / 28,1 / 38,0;
  Alta (VH)   / darkColor              / 30,0 / 40,1 / 50,1
}
\end{verbatim}

\pxConfusionMatrix[title={Matriu de confusi\'o i m\`etriques F1}]
  {Positiu (1), Negatiu (0)}
  {3 (TP), 2 (FP); 5 (FN), 5 (TN)}
  [Precisi\'o = 0.60 \quad|\quad Sensibilitat = 0.375 \quad|\quad $F_1 = 0.4615$]

\pxFuzzyPlot[title={Funcions de pertinen\c{c}a borrosa}, xlabel={Temperatura ($^\circ$C)}]{0, 50}{
  Baixa (VL)  / primaryColor           / 0,1 / 10,1 / 20,0;
  Mitjana (M) / secondaryColor!75!black / 12,0 / 22,1 / 28,1 / 38,0;
  Alta (VH)   / darkColor              / 30,0 / 40,1 / 50,1
}

\end{document}
